\section{交易预案：先选打法，再等形态}

龙头战法在这个题材里只保留两种打法。两种打法都以收盘确认作为前提，盘中情绪、封单和传闻只负责提示观察，不能替代确认。

\begin{exhibitbox}[展品：两种可研究的打法]
\begin{tabularx}{\textwidth}{L{2.3cm}X X}
\toprule
打法 & 候选激活条件 & 盘后撤销条件 \\
\midrule
高度龙头 & 出现项目直接关系主板首板；次日收盘确认二板，且同层有首板补位。 & 二板失败；同层无补位；公司公告澄清关系。 \\
容量核心 & 出现项目直接关系主板股；连续两日保持该层成交额前二，收盘不弱于前收，且未全天一字锁住。 & 收盘弱于前收；成交额跌出该层前二；高度股和容量股同时走弱。 \\
\bottomrule
\end{tabularx}
\end{exhibitbox}

实际复盘顺序如下。

\begin{enumerate}
  \item 先核验公司关系：主板候选必须有公司公告、正式项目名单或同级资料，才能归入项目直接关系层。
  \item 再看梯队：优先看“首板到二板”的连续性，再查新首板是否接上。高度断了，容量股也失去题材溢价基础。
  \item 最后看收盘数据：记录前收、收盘、成交额、连板状态和是否全天一字。收盘弱于前收且梯队转弱，题材观察撤销。
  \item D+2仍看不到第二催化或完整梯队，初始计划作废。D+3到D+5出现第二个官方催化，才检查强化情景。
\end{enumerate}

无创业板权限时，执行顺序更简单：先等项目直接关系的主板节点，再从“高度”与“容量”中选一条线做研究。此节点出现前，掌阅科技只用于观察AI内容扩散，博纳影业和中广天择只用于观察影视扩散。此时最好的预案就是空仓等待，不强行参与。
